\section{Jet extensions by arbitrary monomial order ideals}

\subsection*{Finding M-018: a finite support criterion for detected coefficient jets}
Fix \(r\geq2\), \(\ell\geq0\), and put
\[
 R=r-1,\qquad b=2\ell+2,\qquad
 T=\lceil\sqrt R\rceil\binom{R+\ell}{\ell}.
\]
Let \(A\subset\mathbb Z_{\geq0}^R\) be a nonempty finite
monomial order ideal: it contains every componentwise smaller
multi-index than any of its members. Write
\(D_b=\{\eta:|\eta|\leq b\}\), and let
\(\Gamma\subseteq A+D_b\) be any finite target support.
Then the determinant coefficient projection onto \(\Gamma\)
is dominant at every size
\[
                         n\geq N:=T(1+|A|).                 \tag{M-018.1}
\]
The differential has an elementary unit minor over \(\mathbb Z\).
With the internally padded permanent, the same target projection
is dominant whenever \(n\geq m\geq N\).

In particular, let \(\lambda\vdash nd\) have length at most \(r\),
and define its nested coefficient support
\[
 \Gamma_\lambda=\left\{\gamma\in\mathbb Z_{\geq0}^{r-1}:
       \sum_{i=j}^r\gamma_i\leq\lambda_j\quad(2\leq j\leq r)
                   \right\}.
                                                               \tag{M-018.2}
\]
If \(\Gamma_\lambda\subseteq A+D_b\), (M-018.1) implies
\(K_n(\lambda)=a_\lambda(d[n])\). Under the permanent-size
hypothesis it implies \(q^m_\lambda=K_n=a_\lambda\).
The support condition is finite and explicit; it need not have
an efficient optimization algorithm. M-017 is its special case
where \(A\) is a simplex in an initial segment of the variables.

\begin{proof}
Start with M-013's acyclic graph, with \(T\) inward and
\(T\) outward vertices. Its selected inward-to-outward inverse
entries are pure monomials of every degree one through \(b-1\).
At each inward vertex \(v\), attach a separate tree with vertices
indexed by \(A\), identifying its zero vertex with \(v\).
For every nonzero \(\alpha\in A\), choose any supported index
\(i\); its parent \(\alpha-e_i\) lies in \(A\).
Orient the edge toward the parent and label it by \(z_i\).
Different attached trees have no common nonroot vertices.
There are \(T(|A|-1)\) new vertices, proving the size (M-018.1).

The graph is still acyclic, and no old edge enters a new nonroot
vertex. Consequently an old selected monomial inverse entry
\(z^\eta\), extended at its inward source by \(\alpha\),
becomes exactly \(z^{\alpha+\eta}\). There is one directed
path, not a sum of paths. There is also exactly one path from
the new vertex \(\alpha\) to its own old root \(v\), with
monomial \(z^\alpha\). When \(\alpha=0\), the corresponding
diagonal cofactor is one.

For \(\gamma\in\Gamma\), choose a decomposition
\(\gamma=\alpha+\eta\), \(\alpha\in A\), \(|\eta|\leq b\).
If \(|\eta|\geq2\), remove one variable from \(z^\eta\)
and use that variable as the linear entry variation \(H\).
The remaining degree is between one and \(b-1\), hence is
an old selected inverse entry. Its extension by \(\alpha\)
gives derivative exactly \(z^\gamma\).
If \(|\eta|\leq1\), use the path from \(\alpha\) to its
own old root, varying the reverse entry by \(H=z^\eta\)
when \(|\eta|=1\), and by the homogenizing variable \(x\)
when \(\eta=0\). These again give exactly \(z^\gamma\)
after \(x=1\).

For \(xI-B\), an acyclic path gives its positive monomial
determinant cofactor, as proved in M-013. For \(xI+B\),
the same path gives its positive permanent minor, as proved
in F-021. There are no disjoint nontrivial cycles in either
case. Every target monomial therefore is the derivative of
an elementary source coordinate, with coefficient one and
no other monomial. Distinct target monomials cannot use the
same source coordinate. Their Jacobian minor is the identity.

The degree bound is automatic. A longest new path, of length
\(\max_{\alpha\in A}|\alpha|\), can be attached at the
source of an old degree-\(b-1\) path. Their concatenation
is simple, so \(N\geq\max_{\alpha\in A}|\alpha|+b\).
In particular every target monomial has degree at most \(N\).
Append \(xI\) blocks at larger sizes. Translation and restriction
to the selected directions give \(F(t)=t+O(t^2)\) integrally;
the least homogeneous term proves injective coordinate pullback
after every commutative base change.

For the internally padded permanent, hold one old outward
diagonal entry fixed at \(x\), and move it to position \((1,1)\).
Every nonconstant selected path has distinct endpoints; the
constant cofactor can be taken at an inward diagonal. No selected
direction changes the fixed outward diagonal. Thus multiplication
by its \((n-m)\)-th power preserves all dehomogenized target
coefficients and their dominant projection.

Finally, F-016 places every highest-weight polynomial of type
\(\lambda\) in the algebra generated by the coordinates
(M-018.2). The proved pullback is injective on this entire
polynomial algebra, so no nonzero highest-weight ideal vector
of that type vanishes on the corresponding orbit closure.
This proves all multiplicity assertions, copy by copy.
\end{proof}

\paragraph{A concrete further BIP restriction.}
Every actual obstruction at \(n\geq m^{25}\) satisfies
\[
 \lambda_2\geq1000001\quad\text{or}\quad
 \lambda_3\geq29\quad\text{or}\quad\lambda_4\geq27.
                                                               \tag{M-018.3}
\]
Indeed use \(b=26\), and let \(A\) consist of monomials
\(x_2^a x_3^c\) with
\[
                  a+c\leq999974,\qquad c\leq28.
\]
This order ideal has
\[
 |A|=\sum_{c=0}^{28}(999975-c)=28998869,
                \qquad 1+|A|=28998870<478126755=C_*.
\]
If all three inequalities in (M-018.3) fail, every generator
\(\gamma\in\Gamma_\lambda\) has degree at most one million,
suffix degree from coordinate four at most twenty-six, and
individual coordinate-three exponent at most twenty-eight.
When \(|\gamma|>26\), choose an early submonomial of degree
\(|\gamma|-26\); it lies in \(A\). When \(|\gamma|\leq26\),
use the zero member of \(A\). Thus
\(\Gamma_\lambda\subseteq A+D_{26}\).
The uniform size estimate of M-017, with \(m\geq190\),
gives \(N\leq m^{25}\), and the full determinant multiplicity
excludes an obstruction. Partitions with fewer than four rows
were already excluded by O-006--O-007; zero padding gives the
same argument whenever needed.

\paragraph{Geometric dictionary.}
An order ideal indexes a finite free monomial coefficient module.
The graph constructs an integral morphism between affine spaces,
and its elementary unit Jacobian proves universally injective
coordinate pullback. The multigrading is a torus comodule in the
sense of SGA~3, I, Proposition~4.7.3, pp.~35--36. The vanishing
of the complete ideal multiplicity space uses the comodule
formalism of Proposition~4.7.2, pp.~34--35, and the Schur
operations of SGA~6, VI, Theorem~3.3, pp.~392--393, through
F-016. The decisive data are actual maps of coefficient schemes.
The assertion concerns dominance of a projection; it does not
give exact representations of all polynomials with support
\(\Gamma\), nor a saturation theorem for virtual classes.
